\documentclass[11pt]{article}
\usepackage[margin=1in]{geometry}
\usepackage{amsmath,amssymb,amsthm}
\usepackage{hyperref}
\usepackage{enumitem}

\newtheorem{theorem}{Theorem}
\newtheorem{lemma}[theorem]{Lemma}
\newtheorem{corollary}[theorem]{Corollary}
\theoremstyle{definition}
\newtheorem{definition}[theorem]{Definition}
\theoremstyle{remark}
\newtheorem{remark}[theorem]{Remark}

\usepackage{xurl}
\let\oldus\_
\renewcommand{\_}{\oldus\allowbreak}

\title{A disproof of Conjecture 00000003963\\[2pt]
\large The algebraic connectivity of a hypercube layer is not $\Theta(\log n/n^2)$}
\date{}

\begin{document}
\sloppy
\maketitle

\section{The conjecture}

Conjecture 00000003963 reads:

\begin{quote}
\emph{Definition:} A cube layer is the set of vertices of fixed weight in the hypercube.
\emph{Conjecture:} The algebraic connectivity of the induced graph of a layer, a Johnson graph,
attains its minimum over all layers at $l = n/2$, of order $\Theta(\log n / n^2)$.
\end{quote}
The Chinese version says the same: the algebraic connectivity of the induced graph (Johnson graph) of
layer $\ell$ is smallest at $\ell = n/2$ among all layers, and its order is $\Theta(\log n/n^2)$.

We identify a vertex of the hypercube $Q_n=\{0,1\}^n$ with the subset of $[n]=\{0,\dots,n-1\}$ where it
equals $1$; its weight is the size of the subset. The layer of weight $\ell$ is the set of
$\ell$-subsets of $[n]$.

\paragraph{Definitions used.}
\begin{itemize}[leftmargin=*]
\item \emph{Hypercube} $Q_n$: subsets $s,t\subseteq[n]$ are adjacent iff $|s\,\triangle\,t|=1$
  (they differ in exactly one coordinate).
\item \emph{Johnson graph} $J(n,\ell)$: vertices are the $\ell$-subsets of $[n]$, and $A\sim B$ iff
  $|A\cap B|=\ell-1$. Wikipedia (\url{https://en.wikipedia.org/wiki/Johnson_graph}): ``two vertices are
  adjacent when the intersection of the two vertices (subsets) contains'' $(k-1)$ elements.
\item \emph{Laplacian} $L=D-A$ (Mathlib's \texttt{SimpleGraph.lapMatrix}).
\item \emph{Algebraic connectivity} $a(G)=\lambda_2(L)$. Wikipedia
  (\url{https://en.wikipedia.org/wiki/Algebraic_connectivity}): ``The algebraic connectivity (also known
  as Fiedler value or Fiedler eigenvalue after Miroslav Fiedler) of a graph G is the second-smallest
  eigenvalue (counting multiple eigenvalues separately) of the Laplacian matrix of G.''
\item \emph{Normalized Laplacian} $\mathcal L = I - D^{-1/2}AD^{-1/2}$: entry $1$ on the diagonal
  ($0$ at an isolated vertex), $-1/\sqrt{d_id_j}$ for adjacent $i\ne j$, and $0$ otherwise.
\end{itemize}

\paragraph{Readings.} We refute three readings.
\begin{itemize}[leftmargin=*]
\item[(J)] The appositive ``a Johnson graph'' fixes the layer graph as $J(n,\ell)$, and $a=\lambda_2(L)$.
\item[(N)] The same graph, with $\lambda_2$ of the normalized Laplacian $\mathcal L$.
\item[(I)] The literal reading: the subgraph of $Q_n$ induced on the layer, with $a=\lambda_2(L)$.
\end{itemize}
Since $\ell=n/2$ must be an integer, the order statement is read along even $n$. $f=\Theta(g)$ means
$cg(n)\le f(n)\le Cg(n)$ for all large even $n$, with constants $c,C>0$. Under (J) and (N) we refute
the upper half $f=O(\log n/n^2)$, and under (I) the lower half. Either refutes the conjunction.

\paragraph{Scope.} Under (J), every non-trivial layer has $\lambda_2=n$. So $\ell=n/2$ \emph{is} a
minimiser over the layers $1\le\ell\le n-1$, though not the unique one, and we do not refute the
location clause. The conjecture fails through its order clause. Layers $0$ and $n$ are single
vertices, where $\lambda_2$ is undefined. In Lean, \texttt{lambda2} is set to $0$ when there are fewer
than two eigenvalues, and these layers are excluded from the minimum ($1\le k\le n-1$).

\section{From a quadratic-form bound to $\lambda_2$}

In Lean, $\lambda_2$ of a real symmetric matrix $M$ indexed by $V$ with $|V|\ge2$ is
\texttt{hM.eigenvalues\_0} at index $|V|-2$. Mathlib's list \texttt{eigenvalues\_0} is sorted in
decreasing order (\texttt{eigenvalues\_0\_antitone}), so this is the second-smallest eigenvalue counted
with multiplicity.

\begin{lemma}[\texttt{lambda2\_ge}]\label{lem:var}
Let $M$ be real symmetric with $M\mathbf 1=0$, let $c>0$, and suppose $c\sum_i x_i^2\le x^\top Mx$ for
every $x$ with $\sum_i x_i=0$. If $|V|\ge 2$, then $c\le\lambda_2(M)$.
\end{lemma}
\begin{proof}
Let $(v_j)$ be Mathlib's orthonormal eigenbasis, $Mv_j=\mu_jv_j$. Since $M$ is symmetric and
$M\mathbf 1=0$, we have $\mathbf 1^\top Mx=0$ for all $x$, so $\mu_j\sum_t v_j(t)=0$. If $\mu_j\neq0$,
then $v_j\perp\mathbf 1$ and $c=c\|v_j\|^2\le v_j^\top Mv_j=\mu_j$. So every eigenvalue is $0$ or at
least $c$. Moreover, $\mu_j=0$ forces $\sum_t v_j(t)\ne0$: otherwise $c\le v_j^\top Mv_j=0$. Suppose
$i\ne j$ and $\mu_i=\mu_j=0$. Put $a=\sum v_j$, $b=\sum v_i$ and $x=av_i-bv_j$. Then $\sum x=0$ and
$Mx=0$, so $c(a^2+b^2)=c\|x\|^2\le0$, which contradicts $b\ne0$. Hence at most one index carries an
eigenvalue below $c$. Because the list is sorted in decreasing order, the entries at indices $|V|-2$
and $|V|-1$ come from two different indices. So the entry at $|V|-2$ is at least $c$.
\end{proof}

\section{The Johnson graph: $\lambda_2\ge n$ for all $n$}

For $f$ defined on subsets of $[n]$, put
\[(Df)(C)=\sum_{i\notin C}f(C\cup\{i\}),\qquad (Uf)(E)=\sum_{i\in E}f(E\setminus\{i\}),\qquad
X f(A)=\sum_{j\in A}\sum_{i\notin A}f\big((A\setminus\{j\})\cup\{i\}\big).\]
Write $\mathcal S_m$ for the family of $m$-subsets of $[n]$.

\begin{lemma}\label{lem:alg}
(a) \texttt{up\_dn}, \texttt{dn\_up}: $UDf(A)=|A|f(A)+Xf(A)$ and $DUf(A)=(n-|A|)f(A)+Xf(A)$.
(b) \texttt{adjoint}: $\sum_{E\in\mathcal S_{m+1}}g(E)\,Uh(E)=\sum_{A\in\mathcal S_m}h(A)\,Dg(A)$.
(c) \texttt{sum\_dn}: $\sum_{\mathcal S_m}Dg=(m+1)\sum_{\mathcal S_{m+1}}g$.
\end{lemma}
\begin{proof}
(a) Expand both double sums; the diagonal term ($i=j$) gives $f(A)$ each time. (b) Both sides sum over
pairs $(E,i)$ with $i\in E$, which correspond to pairs $(A,i)$ with $i\notin A$ via $A=E\setminus\{i\}$.
(c) Apply (b) with $h\equiv1$.
\end{proof}

\begin{lemma}[\texttt{key}]\label{lem:key}
If $m+1\le n$ and $\sum_{\mathcal S_{m+1}}f=0$, then
$\sum_{\mathcal S_m}(Df)^2\le m(n-m-1)\sum_{\mathcal S_{m+1}}f^2$.
\end{lemma}
\begin{proof}
We induct on $m$. For $m=0$, $\mathcal S_0=\{\emptyset\}$ and $Df(\emptyset)=\sum_{\mathcal S_1}f=0$ by
Lemma~\ref{lem:alg}(c). For the step, let $f$ live on $\mathcal S_{m+2}$ with sum $0$, and let $z=Df$,
which has sum $0$ by (c). Put $Z=\|z\|^2$ and $F=\|f\|^2$. By (b), $Z=\sum_{\mathcal S_{m+2}}f\,Uz$.
By (b) twice and (a), $\|Uz\|^2=\|Dz\|^2+(n-2(m+1))Z$. The induction hypothesis then gives
$\|Uz\|^2\le (m+1)(n-m-2)Z$. Cauchy--Schwarz gives $Z^2\le F\|Uz\|^2\le(m+1)(n-m-2)FZ$, so
$Z\le(m+1)(n-m-2)F$ (trivially if $Z=0$).
\end{proof}

\begin{theorem}[\texttt{johnson\_quad}, \texttt{algConn\_johnson\_ge}]\label{thm:J}
For $k=m+1\ge1$ and every $x$ on $\mathcal S_k$ with $\sum x=0$: $n\|x\|^2\le x^\top L x$, where $L$ is
the Laplacian of $J(n,k)$. Hence $\lambda_2(L(J(n,k)))\ge n$ for all $n$ and all $1\le k\le n-1$.
\end{theorem}
\begin{proof}
The neighbours of $A$ in $J(n,k)$ are exactly the sets $(A\setminus\{j\})\cup\{i\}$ with $j\in A$ and
$i\notin A$, and this description is injective (\texttt{nbr\_sum}). Hence $\deg A=k(n-k)$
(\texttt{degree\_johnson}) and $\sum_{B\sim A}x_B=Xx(A)$. By Lemma~\ref{lem:alg}(a),
$(Lx)(A)=k(n-k+1)x_A-UDx(A)$, and by (b), $x^\top Lx=k(n-k+1)\|x\|^2-\|Dx\|^2$. Lemma~\ref{lem:key}
gives $\|Dx\|^2\le m(n-m-1)\|x\|^2$, and $(m+1)(n-m)-m(n-m-1)=n$. Finally, $L\mathbf 1=0$ and
$|\mathcal S_k|=\binom nk\ge2$, so Lemma~\ref{lem:var} applies with $c=n$.
\end{proof}

\begin{remark}
The vector $x_A=[0\in A]-k/n$ satisfies $Lx=nx$, so in fact $\lambda_2=n$. Lean proves only the lower
bound, which is all the disproof needs.
\end{remark}

\begin{corollary}[\texttt{normConn\_johnson\_ge}]
For $1\le k\le n-1$, $\lambda_2(\mathcal L(J(n,k)))\ge n/(k(n-k))$; for even $n$ and $k=n/2$ this is $4/n$.
\end{corollary}
\begin{proof}
$J(n,k)$ is $d$-regular with $d=k(n-k)>0$, so $\mathcal L=d^{-1}L$ (\texttt{normLap\_johnson}). Apply
Lemma~\ref{lem:var} with $c=n/d$, using Theorem~\ref{thm:J}.
\end{proof}

\section{The literal reading}

\begin{theorem}[\texttt{algConn\_inducedLayer}]
For all $n,\ell$, the subgraph of $Q_n$ induced on the layer of weight $\ell$ has no edges. Its Laplacian
is $0$, and its $\lambda_2$ is $0$.
\end{theorem}
\begin{proof}
If $|s|=|t|=\ell$, then $|s\setminus t|=|t\setminus s|$, so $|s\triangle t|=2|s\setminus t|\ne1$
(\texttt{inducedLayer\_no\_adj}). Every eigenvalue of the zero matrix is $0$ (\texttt{lambda2\_zero}).
\end{proof}

\section{Asymptotics and the main theorem}

\begin{lemma}[\texttt{exists\_large}]
For every $C\in\mathbb R$ and $N\in\mathbb N$, there is an even $n\ge\max(N,2)$ with $C\log n<4n$.
\end{lemma}
\begin{proof}
Take $n=2(N+M+1)$ with $M\ge C^2$. Since $\log\sqrt n\le\sqrt n-1$, we get $\log n<2\sqrt n$. If
$C\le0$, the claim is clear. Otherwise $C\le\sqrt n$, so $C\log n<2\sqrt n\sqrt n=2n<4n$.
\end{proof}

\begin{theorem}[\texttt{conjecture\_3963\_false}]
All four of the following statements are false:
\begin{enumerate}[leftmargin=*]
\item (J, minimum over layers; \texttt{johnson\_min\_not\_bigO}) there are $C,N$ such that for every even
  $n\ge N$ some layer $1\le k\le n-1$ has $\lambda_2(L(J(n,k)))\le C\log n/n^2$;
\item (J, middle layer; \texttt{johnson\_not\_bigO}) there are $C,N$ with
  $\lambda_2(L(J(n,n/2)))\le C\log n/n^2$ for all even $n\ge N$;
\item (N, middle layer; \texttt{johnson\_norm\_not\_bigO}) the same bound for $\lambda_2(\mathcal L(J(n,n/2)))$;
\item (I, middle layer; \texttt{induced\_not\_bigOmega}) there are $c>0$ and $N$ with
  $c\log n/n^2\le\lambda_2$ of the induced layer graph of weight $n/2$, for all even $n\ge N$.
\end{enumerate}
\end{theorem}
\begin{proof}
(1)--(2): By Theorem~\ref{thm:J}, every non-trivial layer has $\lambda_2\ge n$. With $n$ from the lemma,
$C\log n<4n\le n^3$, so $C\log n/n^2<n$. (3): For $n=2t$, $\lambda_2\ge 2t/t^2$, and
$C\log n\cdot t^2<8t\cdot t^2$ gives $C\log n/n^2<2t/t^2$. (4): $\lambda_2=0<c\log n/n^2$ for $n\ge2$.
\end{proof}

Hence the conjectured order $\Theta(\log n/n^2)$ fails under readings (J), (N) and (I). Under (J) and
(N) the true order is $n$ and $4/n$ respectively; under (I) the value is identically $0$.

\section{Lean formalization}

The file \texttt{lean/Conjecture3963/Basic.lean} (namespace \texttt{C3963}, only import
\texttt{Mathlib}, 564 lines) defines \texttt{hypercube}, \texttt{Layer}, \texttt{johnson},
\texttt{lambda2}, \texttt{algConn} (\texttt{lambda2} of \texttt{lapMatrix}), \texttt{normLap} and
\texttt{inducedLayer}. It proves every statement above under the names given. The main theorem is
\texttt{C3963.conjecture\_3963\_false}; the all-$n$ spectral bound is \texttt{C3963.algConn\_johnson\_ge}.
\texttt{\#print axioms} for both gives exactly \texttt{propext}, \texttt{Classical.choice} and
\texttt{Quot.sound}. There is no \texttt{sorry} and no \texttt{native\_decide}.

\end{document}
